% Copyright (c) 2026 by C. Klukas.
% Licensed under the MIT License. See LICENSE for details.

A paglet is reached by messages whose names and payload types are fixed in
its source. Tools and people who do not read that source need the same
facts in a form a host can read without running the module. This chapter
describes that form, the interface description, and the programs that use
it: the command line \code{paglets-host}, the web UI \code{paglets-web},
and the generator that keeps the documentation and this paper true to the
sources.

\subsection{Interface descriptions}\label{sec:tools-interfaces}

An interface description says which messages a paglet understands, which
types their payloads and replies have, and which arguments it is started
with. It is a JSON document in a Wasm custom section named
\code{paglets.interface}, so it travels inside the module. Custom sections
are part of the module's bytes: the module hash and the ledger's module
signatures (Section~\ref{sec:mesh}) cover the description, and a host reads
it from the bytes it already stores, without instantiating anything.
Listing~\ref{lst:counter-interface} is the description of the counter
example.

% snippet: cpp/examples/counter/interface.json [float,caption={The counter's interface description},label={lst:counter-interface}]
\begin{lstlisting}[style=paglets-json,float,caption={The counter's interface description \src{cpp/examples/counter/interface.json}},label={lst:counter-interface}]
{
  "interface": 1,
  "name": "counter",
  "summary": "A counter whose state is plain C++ objects; memory images move it",
  "launch": {"type": "Start", "example": {"value": 0}},
  "messages": [
    {"name": "increment", "kind": "both", "summary": "Adds to or multiplies the value; replies with the status",
     "request": "Increment", "reply": "Status", "example": {"by": 1, "note": "from the web", "mode": "add"}},
    {"name": "status", "kind": "call", "summary": "The value and the latest notes", "request": "none", "reply": "Status"},
    {"name": "echo", "kind": "call", "summary": "Returns the record unchanged (every field type)", "request": "Echo", "reply": "Echo"},
    {"name": "bloat", "kind": "call", "summary": "Keeps the given amount of memory (image size experiments)",
     "request": "Bloat", "reply": "Status"},
    {"name": "log", "kind": "both", "summary": "Writes a line to the host's log", "request": "string", "reply": "Status"},
    {"name": "schema", "kind": "call", "summary": "The schema descriptor of the messages (JSON text)", "request": "none", "reply": "string"}
  ],
  "types": [
    {"name": "Start", "kind": "record", "fields": [
      {"name": "value", "type": "i64", "doc": "The value to begin with", "default": 0}]},
    {"name": "Increment", "kind": "record", "fields": [
      {"name": "by", "type": "i64", "doc": "Amount (or factor)", "default": 1},
      {"name": "note", "type": "string", "doc": "Kept in the history"},
      {"name": "mode", "type": "Mode", "default": "add"}]},
    {"name": "Bloat", "kind": "record", "fields": [
      {"name": "kilobytes", "type": "u32", "doc": "Memory to keep, in KB", "default": 1024},
      {"name": "fill", "type": "u8", "doc": "Byte value written into it"}]}
  ]
}
\end{lstlisting}
% /snippet

The members are few. \code{interface} is the format version, \code{1}.
\code{launch.type} names the type of the start arguments, or \code{none};
\code{launch.example} fills a launch form. Each entry of \code{messages}
has the name the paglet's router handles, a \code{kind} (\code{call}, a
request with a reply; \code{send}, one-way; or \code{both}; the default is
\code{call}), the \code{request} and \code{reply} types (\code{none} for an
empty payload), one line of \code{summary}, and optionally an
\code{example} payload and \code{danger}, which makes \code{paglets-web} ask
before it sends the message. Names that begin with \code{paglets.} are
reserved. Types are the names of the schema generator
(Section~\ref{sec:model-schemas}): \code{bool}, the integers \code{i8}
to \code{u64}, \code{f32}, \code{f64}, \code{string}, \code{bytes},
\code{list<T>}, \code{optional<T>}, \code{any} for a free JSON value, and
the records and enums listed under \code{types}. Records encode as
MessagePack maps keyed by field name and enums as the enumerator's name,
which are the encodings of the generated codecs, so a description matches
the C++ types of the schema header (Section~\ref{sec:model}).

The description is not the schema. The schema header is the source of the
types; the description adds what the types do not say: the message names,
their kinds, summaries, examples, and per field a \code{doc} and a
\code{default}. The counter's file therefore defines only \code{Start},
\code{Increment} and \code{Bloat}, the records whose fields it documents.
\code{Status}, \code{Echo} and the enum \code{Mode} come from the schema
when the module is built: \code{paglets\_add\_module} takes
\code{INTERFACE counter/interface.json} next to \code{SCHEMA\_HEADER}, and
after linking it runs the interface tool, \code{paglets-interface}, on the
module (Figure~\ref{fig:interface-flow}). The tool merges the records and enums of the generated
\code{counter.schema.json} into \code{types} (a type the description
defines itself takes precedence), checks the result and writes the
section. A description that names an unknown type, repeats a message or
uses another kind fails the build with the reason. A module whose schema
namespace holds service contracts (Section~\ref{sec:mesh}) needs no file:
the tool makes a description with one \code{call} message per operation.
Because embedding changes the module's bytes, a module is signed and
trusted after it was built, as any module is.

\begin{figure}[htbp]
\centering
\begin{tikzpicture}[every node/.append style={font=\footnotesize},
                    box/.append style={inner xsep=3pt}, plain/.append style={inner xsep=3pt},
                    state/.append style={inner xsep=4pt}]
  \node[plain] (json) {\code{interface.json}\\messages, kinds, docs};
  \node[plain, below=3mm of json] (schema) {\code{counter.schema.json}\\records, enums};
  \coordinate (mid) at ($(json.east)!0.5!(schema.east)$);
  \node[box, right=5mm of mid] (tool) {\code{paglets-interface}\\merge, check, embed};
  \node[state, right=5mm of tool] (module) {\code{counter.wasm}\\\code{paglets.interface}};
  \node[plain, above right=1mm and 5mm of module] (cli) {\code{paglets-host}\\\code{module inspect}\\\code{remote describe}};
  \node[plain, below right=1mm and 5mm of module] (web) {\code{paglets-web}\\forms, launch,\\shell commands};
  \draw[arrow] (json) -- (tool);
  \draw[arrow] (schema) -- (tool);
  \draw[arrow] (tool) -- (module);
  \draw[arrow] (module) -- (cli);
  \draw[arrow] (module) -- (web);
\end{tikzpicture}
\caption{From the description and the schema's types (written by the schema generator) into the module at build time, and out of the module into the tools. The section is covered by the module hash and the ledger's module signatures. Hosts read it from the stored bytes; the paglet never runs for it.}
\label{fig:interface-flow}
\end{figure}

In the host tree the description is a small API of the \code{wire} library
over the module bytes, which the interface tool links as well
(Listing~\ref{lst:interface-api}). \code{of\_module} finds and parses the
section; \code{validate} is the check the build runs; \code{from\_schema}
builds the generated description of a contract guest. The same code
answers the \code{describe} request of a remote session, also for the
native system paglets, whose descriptions are built from their service
contracts.

% snippet: cpp/host/include/paglets/wire/interface.hpp#api [float,caption={The host's view of interface descriptions},label={lst:interface-api}]
\begin{lstlisting}[style=paglets-cpp,float,caption={The host's view of interface descriptions \src{cpp/host/include/paglets/wire/interface.hpp}},label={lst:interface-api}]
// Checks a description: the format version, message names and kinds, and
// that every type it names is a built-in type or one of its `types`.
std::expected<void, std::string> validate(const Json& description);

// Parses and validates a description given as JSON text.
std::expected<Json, std::string> parse(std::string_view text);

// Adds the records and enums of a schema descriptor (wire/schema.hpp,
// <name>.schema.json) that the description does not define itself.
Json merge_schema_types(Json description, const Json& schema);

// A description made from a schema descriptor's services: one `call`
// message per operation.
Json from_schema(const Json& schema, std::string_view name);

// The interface of a module's bytes: the custom section, parsed; nullopt
// when the module has none.
std::expected<std::optional<Json>, std::string> of_module(std::span<const std::uint8_t> wasm);
\end{lstlisting}
% /snippet

The build checks that every type the description names exists, in the
file or in the schema, and nothing against the paglet's router: a message
the file lists but no handler takes is found when it is sent, by the
\code{unknown\_message} error reply, not by the build. The router keeps
its handlers in a map by name but exposes neither the map nor the names,
so a test cannot compare the two directly; it could only send each
described message and watch for \code{unknown\_message}. The end-to-end
test of \code{paglets-web} launches the counter with a start value, reads
its description and exercises three of its messages (a send and two
calls); beyond that, a description written by hand is as true as its
author keeps it.

\subsection{The command line}\label{sec:tools-cli}

\code{paglets-host} is one program with four groups of commands
(Table~\ref{tab:cli-groups}). The first group runs paglets on one machine
without a mesh, and so without keys or a ledger. The second manages keys
and the signed ledger of a mesh, and inspects module files, all offline.
\code{serve}
runs a host. \code{remote} talks to a running host.

\begin{table}[htbp]
\centering\small
\begin{tabular}{@{}lll@{}}
\toprule
Group & Commands & Works on \\
\midrule
local & \code{run}, \code{list}, \code{call} & a module file or a state directory here \\
records & \code{keys}, \code{mesh create}, \code{ledger}, \code{module inspect} & key files, a ledger copy or a module file \\
host & \code{serve} & this machine, as a host of a mesh \\
sessions & \code{remote \emph{command}} & a live host, over an end-to-end channel \\
\bottomrule
\end{tabular}
\caption{The command groups of \code{paglets-host}. \code{paglets-worker} is started by \code{paglets-host} for the sandboxed instances and is not meant to be started by hand (its one option, \code{--check-sandbox}, serves the sandbox test).}
\label{tab:cli-groups}
\end{table}

\code{run} creates one paglet from a module, sends it the \code{--call}
requests in order, prints one line per reply, waits until the paglet is
idle and disposes of it, unless \code{--keep} is given. With
\code{--state-dir} the paglets outlive the process and resume from their
last memory image, and \code{list} and \code{call} work on that directory.
Message bodies and arguments are JSON on the command line;
\code{paglets-host} converts them to MessagePack before they reach the
paglet and prints replies as JSON (\code{ok} for an empty reply,
\code{error: \emph{code}} for an error code). \code{\{"\$bytes": "00ff"\}}
is a binary string for \code{bytes} fields.

\code{remote} opens an end-to-end channel to the host named by
\code{--connect}, as the admin or owner whose key file is \code{--key}; the
role comes from the key file. The channel is the one of
Section~\ref{sec:communication-channel}, with the admin's or owner's key
as the client's static key. The client accepts the host only if its key is
enrolled in the ledger copy \code{--ledger}, or is the key given as
\code{--host-key}; the handshake proves that key, so \code{remote} checks
no TLS certificate and has no CA option (hosts check each other's
certificates only when a CA file is given,
Section~\ref{sec:communication-channel}). Within the session,
\code{launch} starts a paglet of the owner from a module file that travels
along, or from the hash of a module the host has, with a passport the
client signs locally; \code{call}, \code{send} and \code{describe} reach a
paglet; \code{dispatch} moves it and waits for the outcome; \code{locate},
\code{pin}, \code{dispose}, \code{terminate}, \code{hosts},
\code{landscape}, \code{modules} and \code{zombies} are among the others;
\code{status}, the ledger's \code{push} and \code{pull}, \code{pins} and
\code{unpin}, \code{slots}, \code{requests} with \code{approve} and
\code{deny}, \code{audit} and \code{push-module} complete the set. An
owner sees and acts on their own paglets, an admin on all of them; the
host checks every request against its ledger: the session's role, the
paglet's owner, and for a launch the passport. A pin set from a session
needs no policy rule; the rule of Section~\ref{sec:communication-pins}
governs pins that paglets request. Listing~\ref{lst:cli-session} shows a session: the counter starts
with a start value and answers the call with the \code{Status} record of
its description.

\begin{lstlisting}[style=paglets-bash,float,caption={A counter on a mesh host from the command line. The reply is the one the counter gives for these arguments (recorded with \code{paglets-host run}); \code{remote describe} prints the module's description as text: Listing~\ref{lst:counter-interface} with the schema's types \code{Mode}, \code{Status} and \code{Echo} merged in.},label={lst:cli-session}]
H=build/linux-gcc16/host/paglets-host
S="--connect https://lab-1:7443 --key olga.key --ledger ledger"
$H remote launch $S build/linux-gcc16/guests/counter.wasm \
    --args '{"value": 200}'                      # prints the paglet ID
$H remote describe $S <paglet-id>
$H remote call $S <paglet-id> increment '{"by": 5, "note": "first", "mode": "add"}'
{"value":205,"history_size":1,"recent_notes":["first"],"average_step":5,"memory_pages":2}
$H remote dispatch $S <paglet-id> lab-2            # moved to lab-2 <key>
\end{lstlisting}

\subsection{The web UI}\label{sec:tools-web}

\code{paglets-web} shows a mesh in the browser. It is a mesh client like
\code{remote}: it opens sessions with the hosts as an admin or owner,
with the same \code{--connect}, \code{--key} and \code{--ledger} options,
and serves its pages itself on this computer, from its own HTTP server
compiled into the program. Nothing is installed on the hosts; they run
\code{paglets-host serve} as before. Listing~\ref{lst:web-app} is the core
of the program.

% snippet: cpp/web/src/app.hpp#app [float,caption={The application behind \code{paglets-web}: a poller and the JSON API},label={lst:web-app}]
\begin{lstlisting}[style=paglets-cpp,float,caption={The application behind \code{paglets-web}: a poller and the JSON API \src{cpp/web/src/app.hpp}},label={lst:web-app}]
class App {
public:
    App(Options options, client::Identity identity);
    ~App();

    // Polls the mesh until stop(): hosts and landscape from the entry
    // host, status from every host it can reach.
    void run_poller();
    void stop();
    // One refresh now (tests, and the first page).
    void refresh();

    // The HTTP routes below /api/ (the guard has checked the request).
    web::Response api(const web::Request& request);

    // The current model as JSON text.
    std::string model() const;
\end{lstlisting}
% /snippet

Everything the pages show comes from \code{model()}, one JSON text that
the poller keeps fresh: it asks the entry host for its registry of hosts
and for \code{mesh-info}'s landscape, then every host it can reach for its
status. Every action is a route of \code{api()}; \code{Options} holds the
command-line options and \code{client::Identity} the opened key and ledger
copy. The HTTP server, the token guard and the page's assets are outside
the class.

The welcome page lists the hosts with their state and counts; a host page
shows its paglets, modules and pins. A paglet page is built from the
interface description: one form per message, with fields from the request
type, a JSON tab for the same value as text, \emph{Call} and \emph{Send},
and a confirmation for messages marked \code{danger}. The actions are
those of the CLI (dispatch, pin, locate, dispose, terminate); the modules
page launches from the \code{launch} type; the zombies page lists the
groups of Section~\ref{sec:zombies}. A paglet without a description gets
a name-and-JSON form.

Next to every view and action stands the \code{paglets-host remote}
command that does the same in a shell, kept current while a form is
edited, and every result shows the command for exactly what was just
done. The commands are rendered by the client library that both programs
share (\code{client::remote\_command}), with the key file, ledger directory
and passphrase file \code{paglets-web} was started with. A reader who
works in the browser thus sees the command for every step and can move to
the shell with it.

\code{paglets-web} acts with the user's key and guards its API
accordingly: it listens on \code{127.0.0.1} by default (\code{::1} and
\code{localhost} are allowed too; any other address needs \code{--listen}
with \code{--allow-remote}), requires a random per-start token (handed to
the browser in the opened link and kept as an \code{HttpOnly},
\code{SameSite=Strict} cookie, or sent as a bearer header by \code{curl}),
requires a \code{Host} header that names its own address, which defeats
DNS rebinding, and refuses state-changing POSTs from another origin. The
hosts trust none of this: they check every request from \code{paglets-web}
as they check one from \code{remote}, by the key of the session.

\subsection{The documentation generator}\label{sec:tools-generators}

The build has two generators. The schema generator of
Section~\ref{sec:model-schemas} runs in the build and writes, next to the
guest's codecs, the JSON descriptor \code{\emph{name}.schema.json}; it is
what the interface tool merges, so the types in a description and the
types in the paglet's code have one origin.

The documentation generator, \code{paglets-docs-sync}, keeps the
documentation and this paper true to the sources; it is run after a
change to marked code, and the test suite checks that its output is
current. Marked regions of the sources (\code{// --8<-- [start:\emph{name}]}
to \code{[end:\emph{name}]}) are written into the Markdown pages and into
the \code{.tex} chapters between comment markers; the generated listings
of this paper are such regions, and their captions end with the source
path; a caption without one (Listing~\ref{lst:cli-session}) marks a
listing typed by hand. Test results come the same way: a test that reports
a value writes it to \code{docs/results/\emph{topic}.json} when
\code{PAGLETS\_RESULTS\_DIR} is set, and the generator renders the file as
a table. \code{cpp/tools/update\_results.sh} runs the example tests in the
host process and in sandboxed workers, the protocol checker of
Section~\ref{sec:zombies}, and then the generator; \code{mesh\_bench.sh}
records the live benchmarks of Section~\ref{sec:evaluation} on real
machines. The \code{docs\_sync} test of \code{ctest} fails while a page or
a chapter is not current, so a code change that invalidates an excerpt
fails the test suite, and the paper's workflow, which runs the same check
before it builds the PDF, until the text is regenerated. An excerpt in the
published PDF therefore equals the source of the commit the PDF was built
from.
